\section{Lean formalization: completed kernel and exact remaining targets}
\label{sec:formalization}

\subsection{The distinction between three kinds of verification}

There are three separate claims in this package.
First, the paper supplies human-readable arithmetic proofs whose programs
and induction invariants are primitive recursive and quantifier-free.
Second, selected general algorithms and finite-table assertions are proved
in ordinary Lean. Third, one can encode the arithmetic proofs in the
first-order object theory implemented by ProveIt.
The first two are delivered; the full third development is a detailed
implementation plan, not a completed proof term.

An ordinary theorem of Lean about natural numbers is not automatically a
derivation in iPRA. In particular, native Lean permits induction on
propositions with quantifiers, even if the proof never invokes
\code{Classical}. The embedded target must instead construct a term of
\[
 \code{Prf PRAAx}\ \Gamma\ \varphi,
 \qquad\text{written }\Gamma\vdash[\code{PRAAx}]\varphi.
 \tag{7.1}\label{eq:7-derivation}
\]
ProveIt's \code{pra\_iff\_native} translates a derivation that already
exists in one of its embedded calculi. It is not a reflection rule that
turns every ambient Lean theorem into \eqref{eq:7-derivation}.

\subsection{Repository foundations inspected at a fixed revision}

The inspected ProveIt revision is
\begin{center}
\code{27da047166883e5010e115968d3a40f9e91bfb3d}.
\end{center}
The repository was read through the selected GitHub connector. No upstream
files were changed, and the entire repository was not rebuilt.
The following are the exact relevant files under
\path{Logic/PrimitiveRecursiveArithmetic/ConciseIPRA/Lean/ConciseIPRA/}:

\begin{longtable}{@{}p{0.30\textwidth}p{0.64\textwidth}@{}}
\toprule
File & Role in the intended proof\\
\midrule
\endhead
\code{Syntax.lean} & \code{Fm} and the syntactic predicate \code{Fm.QF};
both quantifiers are excluded from the induction formula.\\
\code{Deduction.lean} & Embedded intuitionistic natural deduction,
equality, substitution, weakening and proof composition.\\
\code{PR.lean} & Function-symbol grammar \code{PR}, terms \code{PTm},
axioms \code{PRAAx}, and the derived proof constructor \code{PRA.qfInd}.\\
\code{PRAToScalar.lean} & Compilation of a specified primitive recursive
definition into the alternative iterator language.\\
\code{TeqReflect.lean} & Reflection of quantifier-free formulas into
characteristic terms and the corresponding induction.\\
\code{BiInterpretation.lean} & Translation between the two object theories.\\
\code{NativeBiInterpretation.lean} & Translation between already
constructed embedded derivations, including \code{pra\_iff\_native}.\\
\bottomrule
\end{longtable}

The existing Diophantine-geometry package supplies written bounded-height,
resultant and square-area arguments and a native arithmetic kernel.
Its earlier finite-generation interface assumes an already supplied
primitive recursive reducing selector. The present elliptic development
derives that selector from the actual representative list and its covering
proof. The complete rational-halving algorithm is a further explicit
piece of this derivation. The square-area proof here follows the especially
self-contained version in the repository's Faltings package, with its
exact fixed-fuel invariant. These dependencies and the audit scope are
recorded in \code{metadata/repository\_audit.txt}.

\subsection{The delivered native Lean definitions}

The generic point type is an arbitrary \code{Type}; its laws are supplied
in a structure. The central interface in \code{DescentCore.lean} contains:
\begin{lstlisting}
structure DescentData (alpha : Type u) where
  add : alpha -> alpha -> alpha
  height : alpha -> Nat
  cutoff : Nat
  repCount : Nat
  baseCount : Nat
  rep : Fin repCount -> alpha
  base : Fin baseCount -> alpha
  repIndex : alpha -> Fin repCount
  baseIndex : alpha -> Fin baseCount
  child : alpha -> alpha
  baseComplete :
    forall p, height p <= cutoff -> base (baseIndex p) = p
  decompose :
    forall p, cutoff < height p ->
      add (add (child p) (child p)) (rep (repIndex p)) = p
  contract :
    forall p, cutoff < height p ->
      2 * height (child p) <= height p + cutoff
\end{lstlisting}
This listing transliterates the actual Unicode Lean source for readability.
There is no declaration asserting that every elliptic curve has such a
structure.

The inductive datatype \code{Word r b} has base and representative atoms,
addition, and a compact doubling constructor. \code{Word.eval} interprets
it in the given operation; \code{Word.size} counts constructors;
\code{baseCoeff} and \code{repCoeff} compute its nonnegative coefficients.
The definitions \code{next} and \code{run} freeze inside the base region.
\code{build} recursively constructs the word with numerical fuel, and
\code{coordinates p = build (height p) p}.

The completed theorems include
\[
\begin{split}
 h(\code{run}(n,P))&\le C+(h(P)\dotminus n),\\
 h(\code{run}(h(P),P))&\le C,\\
 \ev(\code{coordinates}(P))&=P,\\
 \operatorname{size}(\code{coordinates}(P))&\le3h(P)+1.
\end{split}
\tag{7.2}
\]
The proof is independent of the group laws. The stronger logarithmic
fuel bound in Theorem~\ref{thm:abstract} is proved in the paper and used
in the arithmetic reference algorithm; it is not advertised as a theorem
already contained in the native kernel.

\code{FiniteSearch.lean} defines a structurally recursive \code{first?},
proves soundness and completeness on an explicit finite list, then supplies
total selection with a default and bounded-natural search.
\code{E1FiniteTable.lean} defines four labels and checks the table,
commutativity, associativity, exponent two, and a two-generator coding.
It does not infer that these labels exhaust the rational curve.
\code{DescentExample.lean} executes the kernel on a simple numerical
instance; \code{Audit.lean} prints the actual axiom dependencies.

The sixth source, \code{EllipticSpec.lean}, supplies concrete native
definitions of integer, pair, rational and point codes; curve membership;
chord addition, negation and doubling; $H$ and bounded-search $h$; the
halving quartic; the covering predicate; finite division candidates; and
generator-word syntax. Its remaining arithmetic obligations are
definitions of propositions, including \code{HalvesCorrectTarget} and
\code{E1ClassificationTarget}. They are not asserted axioms or proved
theorems. Native \code{Rat} is used as an implementation specification;
its standard-library proof fields have dependencies from Lean's fuller
foundations, including classical choice. The first-order arithmetic bridge
must replace these with the explicit numerical coding proofs. This does
not affect the separate descent/search axiom audit.

Compilation used Lean 4.32.0, commit
\code{8c9756b28d64dab099da31a4c09229a9e6a2ef35}.
The main descent and finite-search theorems report only
\code{propext} and \code{Quot.sound}; the finite table's
\code{two\_generators} theorem reports no axioms.
There are no \code{sorry} placeholders or new axiom declarations in the
compiled files, and no reported \code{Classical.choice} or \code{sorryAx}.
This is not a claim that the entire arithmetic paper is formally proved.

\subsection{All data needed for the arithmetic bridge}

Use the concrete encodings of Section~\ref{sec:foundations}.
The intended embedded development has fixed symbols for:

\begin{enumerate}
\item signed integer operations, positive division, gcd with Bezout
coefficients, rational normalization and ordered comparisons;
\item pair/list operations, bounded scans and finite polynomial evaluation;
\item \code{valid(A,B,P)}, \code{neg}, \code{add}, \code{twice},
and the projective abscissa, with formulas \eqref{eq:add};
\item $H,h,\mathcal L_E$, the matrix $M_{A,B}$, its determinant,
adjugate multipliers, normalized $L$, $T_R$ and $C$;
\item rational-root enumeration, rational-square testing and
$\Halves_E$, including the exceptional target $O$;
\item the finite-list selector, stopped descent state, compact word,
coefficient vector, and evaluation;
\item the primitive triangle predicate, normalized triangle,
Pythagorean parameters, square-area map and its iterate;
\item the four-point classification code, scale maps and finite-product
coordinate maps.
\end{enumerate}

These definitions are fully specified by formulas and bounded loops in
the paper. Integer square root may use an exhaustive scan for the object
proof even when Python uses a faster implementation.
For words, code the constructors by a finite tag and the pairing function.
Choose each constructor's code larger than its subword codes.
Evaluation and structural tests can be implemented by course-of-values
primitive recursion, storing the computed values for all smaller codes
in a finite list. Thus the native inductive datatype is not assumed as an
extra first-order sort.

Define the object-language cover formula by
\[
 \forall p\bigl(\operatorname{valid}(A,B,p)=1\to
   \exists i\,\exists q\,
   [\,i<|R|\wedge\operatorname{valid}(A,B,q)=1
    \wedge p=\operatorname{add}(\operatorname{twice}(q),R_i)\,]\bigr).
 \tag{7.3}\label{eq:objectcover}
\]
Bounded list membership in other formulas means equality of a primitive
recursive Boolean fold to $1$. In particular the successful selector
predicate and all induction invariants are quantifier-free.

The main elliptic target is a derivation under
\eqref{eq:objectcover}, nonsingularity and validity of the given list.
Abbreviate the computed terms by
$W=\operatorname{coordinates}(A,B,R,p)$ and $G=G_E(R)$. The target is
\[
 \operatorname{valid}(A,B,p)=1
 \ \longrightarrow\
 \bigl[\operatorname{wordValid}(W,G)=1\ \wedge\ \ev(G,W)=p\bigr].
 \tag{7.4}
\]
Universal closures over the input codes are introduced after the proof.
For the unconditional target replace the hypotheses by
$\operatorname{valid}(-1,0,p)=1$ and conclude that the four-case classifier
decodes to $p$. This gives the bit-coordinate theorem directly.

\subsection{A dependency-ordered implementation path}

\begin{longtable}{@{}p{0.21\textwidth}p{0.73\textwidth}@{}}
\toprule
Module to implement & Exact mathematical obligation and proof route\\
\midrule
\endhead
\code{Codes} & Encode integers, pairs, lists and rationals; prove projection
identities, normalization idempotence, and reduced-fraction uniqueness.\\
\code{Arithmetic} & Order, divisibility and gcd; prove the Euclidean
coefficient invariant and the three square-splitting lemmas.\\
\code{FiniteLoops} & Compile scans and folds into \code{PR} terms;
prove their Boolean prefix invariants with \code{PRA.qfInd}.\\
\code{EllipticOps} & Define the total guarded chord formulas, closure,
inverse translation and projective duplication. General associativity is
needed only for the coefficient-vector group presentation.\\
\code{PointBoxes} & Prove the square/cube denominator lemma and the
explicit $(a,b,d)$ enumeration bounds.\\
\code{PolyCertificates} & Fixed-size determinant and adjugate identities,
the explicit discriminant polynomial, normalization and coefficient norms.\\
\code{HeightBounds} & Prove image-gcd divisibility, duplication lower
bound, translation upper bound including the exceptional cases, and
binary-height contraction.\\
\code{Halving} & Prove the rational-root divisibility bounds, square test,
quartic completeness and the exact theorem
$Q\in\Halves_E(T)\leftrightarrow Q\in E(\Q)\wedge2Q=T$.\\
\code{CoverSelector} & From the supplied cover, prove success of the
finite scan through $R_i,\Halves_E(P-R_i)$ and prove its contraction.\\
\code{DescentIPRA} & Prove the forward-state invariant by quantifier-free
induction; construct the final word or coefficient vector and its bound.\\
\code{SquareArea} & Normalize triangles, derive Pythagorean parameters,
prove \eqref{eq:fermat-map}, and derive the contradiction at
the initial hypotenuse using \eqref{eq:fermat-invariant}.\\
\code{E1AndProducts} & Clear denominators in \eqref{eq:triangle-map};
prove all points lie in the four-element table; define inverse bit maps,
scale maps and product folds.\\
\bottomrule
\end{longtable}

Each row has explicit premises and a mathematical proof in this paper.
For example, the target named \code{Halves\_correct} is a missing
formal arithmetic lemma, not a theorem already present in the compiled
kernel. Once it and the height lemmas are constructed, a target
\code{Cover\_implies\_search\_success} is proved by existential elimination
from \eqref{eq:objectcover}, finite-scan completeness, and inverse
translation. No unbounded search or choice axiom is required.

\subsection{The critical induction translation}

Native Lean's proof of word correctness inducts on fuel while generalizing
over the point. Its induction proposition therefore has a universal
quantifier. The safe embedded implementation instead fixes the initial
point $P$ and numerically iterates one forward state. At stage $j$, a
coarse invariant is
\[
\begin{gathered}
 \operatorname{valid}(Q_j)=1,\quad k_j\le j,\quad
 h(Q_j)\le C+(h(P)\dotminus j),\\
 P=2^{k_j}Q_j+\sum_l a_{j,l}R_l .
\end{gathered}
\tag{7.5}\label{eq:embedded-invariant}
\]
The symbol for the finite sum, all state entries, and validity are fixed
primitive recursive terms. Hence this is a quantifier-free formula with
$P,R,C$ as parameters. For the operation-only version maintain equality
between $P$ and an accumulated context applied to $Q_j$; the context is
a coded compact word with a hole. Its update substitutes
$(Q_{j+1}\oplus Q_{j+1})\oplus R_i$ for the hole.
One can also retain the entire finite trace and verify its equations by a
Boolean fold.

For the optimized fuel, replace the height line by
$2^j(h(Q_j)\dotminus C)\le h(P)\dotminus C$.
For the Fermat proof use exactly \eqref{eq:fermat-invariant}.
These choices resolve the weak-induction issue explicitly rather than
assuming native Lean induction is an admissible object rule.

\subsection{Completion criteria}

A complete embedded formalization should deliver actual \code{Prf PRAAx}
values for the relative elliptic theorem and unconditional classification,
with the primitive recursive symbols and quantifier-free side conditions
defined in the repository. Rebuild those modules, inspect every added axiom
and placeholder, and only then apply the existing interpretation bridges.
The supplied native files and executable arithmetic are useful intermediate
artifacts; they do not meet this final criterion on their own.

Completing the missing finite algebra and coding is a substantial but
specified formalization project. Removing the representative-cover premise
for all elliptic curves is a different mathematical problem and is not
presented as a routine remaining Lean task.
